\speciestitle{Methyl chloride}{CH\cs3Cl}{CH3Cl}%
\begin{description}
\item[Useful range:] 147\,--\,4.6\,hPa
\item[Contact:] Michelle Santee, \email{Michelle.L.Santee@jpl.nasa.gov}
\end{description}
\label{sec:StandardCH3Cl}

\subsection*{Introduction}

The v2.2 MLS ClO measurements were characterized by a substantial
(\texttildelow0.1\,--\,0.4\,ppbv) negative bias at retrieval levels below (i.e.,
pressures larger than) 22\,hPa.  \citet{SanteeEtal08clo} suggested that
contamination from an interfering species such as CH\cs3Cl, which has
lines in two wing channels of the \mbox{640-GHz} band used to measure
ClO, could have given rise to the bias; they showed results from early
v3 algorithms in which CH\cs3Cl was also retrieved that demonstrated
significant reduction in the bias in lower stratospheric ClO\@.  Further
refinements in the \thisv algorithms yielded not only an improved ClO
product, but also a reliable retrieval of CH\cs3Cl.

As for ClO, the standard CH\cs3Cl product is derived from radiances
measured by the radiometer centered near 640\,GHz. The MLS \thisv CH\cs3Cl
data are scientifically useful over the range 147 to 4.6\,hPa.  A summary
of the precision and resolution (vertical and horizontal) of the \thisv
CH\cs3Cl measurements as a function of altitude is given in
Table~\ref{tab:CH3ClSummary}.  More details on the quality of the MLS
\thisv CH\cs3Cl measurements are given below.

\subsection*{Resolution}
\onestandardavk{CH3Cl}{CH\cs3Cl}%

The resolution of the retrieved data can be described using ``averaging
kernels'' \citep[e.g.,][]{Rodgers00}; the two-dimensional nature of the
MLS data processing system means that the kernels describe both vertical
and horizontal resolution.  Smoothing, imposed on the retrieval system
in both the vertical and horizontal directions to enhance retrieval
stability and precision, degrades the inherent resolution of the
measurements.  Consequently, the vertical resolution of the \thisv
CH\cs3Cl data, as determined from the full width at half maximum of the
rows of the averaging kernel matrix shown in Figure~\ref{fig:AvkCH3Cl},
is \texttildelow4\,--\,6\,km in most of the lower stratosphere, degrading to
7\,--\,10\,km at and above (i.e., at pressures less than) 14\,hPa.  Note
that there is overlap in the averaging kernels for the 100 and 147\,hPa
retrieval surfaces, indicating that the 147\,hPa retrieval does not
provide as much independent information as is given by retrievals at
higher altitudes.  Figure~\ref{fig:AvkCH3Cl} also shows horizontal
averaging kernels, from which the along-track horizontal resolution is
determined to be \texttildelow600\,km at 147\,hPa, \texttildelow450\,--\,500\,km from
100 to 22\,hPa, and \texttildelow550\,--\,850\,km at and above 15\,hPa.  The
cross-track resolution, set by the width of the field of view of the
\mbox{640-GHz} radiometer, is \texttildelow3\,km.  The along-track separation
between adjacent retrieved profiles is 1.5\degsym\ great circle angle
(\texttildelow165\,km), whereas the longitudinal separation of MLS
measurements, set by the Aura orbit, is 10\degsym--\,20\degsym\ over low
and middle latitudes, with much finer sampling in the polar regions.

\subsection*{Precision}

\begin{figure}
\centering\dontincludegraphics[width=12.0cm]{figures/MLSCH3Clascdsc-200502}
\caption{(left) Ensemble mean profiles for ascending (red) and
  descending (blue) orbit matching pairs of MLS \thisv CH\cs3Cl in the
  latitude range 50\degsym S\,--\,50\degsym N averaged over 30 days in April
  2008.  (right) Mean difference profiles between ascending and
  descending orbits (cyan), the standard deviation about the mean
  difference (orange), and the root sum square of the precisions
  calculated by the retrieval algorithm (magenta).  SD values are scaled
  by 1/$\sqrt{2}$; thus the observed SD represents the statistical
  repeatability of the MLS measurements, and the expected SD represents
  the theoretical \mbox{1-$\sigma$} precision for a single profile.  See
  \citet{LambertEtal07} for details.  }
\label{fig:CH3Clascdsc}
\end{figure}

The precision of the MLS CH\cs3Cl data is estimated empirically by
computing the standard deviation of the differences between matched
measurement points at the intersections of the ascending (day) and
descending (night) sides of the orbit.  Observed scatter, representing
the statistical repeatability of the measurements, is 100\,pptv or less
throughout the vertical domain.  This estimate reflects the precision of
a single profile; in most cases precision can be improved by averaging,
with the precision of an average of $N$ profiles being 1/$\sqrt{N}$
times the precision of an individual profile (note that this is not the
case for averages of successive along-track profiles, which are not
completely independent because of horizontal smearing).  The theoretical
precision reported by the Level~2 data processing system exceeds the
observationally-determined precision throughout the vertical range,
indicating that the smoothing applied to stabilize the retrieval and
improve the precision has a nonnegligible influence.  Because the
theoretical precisions take into account occasional variations in
instrument performance, the best estimate of the precision of an
individual data point is the value quoted for that point in the L2GP
files, but it should be borne in mind that this approach slightly
overestimates the actual measurement noise.

\subsection*{Range}

Although CH\cs3Cl is retrieved (and reported in the L2GP files) over the
range 147 to 0.001\,hPa, on the basis of the drop off in precision and
resolution, the reduction in independent information contributed by the
measurements, and the results of simulations using synthetic data as
input radiances to test the closure of the retrieval system, the data
are not deemed reliable at retrieval levels above (i.e., pressures lower
than) 4.6\,hPa.  Despite the overlap in the averaging kernels for the
147 and 100\,hPa surfaces (Figure~\ref{fig:AvkCH3Cl}), maps at 147\,hPa
display substantial features not seen at 100\,hPa (not shown) that are
believed to represent real atmospheric variations.  Thus we recommend
that the \thisv CH\cs3Cl data may be used for scientific studies between
147 and 4.6\,hPa, inclusive.

\subsection*{Accuracy}

The impact of various sources of systematic uncertainty has not yet been
quantified for CH\cs3Cl as it has for most other MLS products.  This
work is planned as part of a dedicated validation exercise for the \thisv
CH\cs3Cl data.

\subsection*{Review of comparisons with other datasets}

Detailed comparisons with correlative data sets have not yet been
undertaken.  This work is planned as part of a dedicated validation
exercise for the \thisv CH\cs3Cl data.

\subsection*{Data screening}
\begin{description}
\item[Pressure range:]  \screeningrule{147\,--\,4.6\,hPa}

  \standardrangestatement
  
  \standardprecisionrule
  
\item[Status flag:] \screeningrule{Only use profiles for which the
    `Status' field is zero.}  We recommend that all profiles with
  nonzero values of `Status' be discarded, because of the potential
  impact of cloud artifacts at lower levels.  See
  Section~\ref{sec:StatusQuality} for more information on the
  interpretation of the \texttt{Status} field.

\item[Clouds:] \cloudsokrule

  Nonzero but even values of \texttt{Status} indicate that the profile
  has been marked as questionable, usually because the measurements may
  have been affected by the presence of thick clouds.  Globally fewer
  than \texttildelow1\,--\,2\% of CH\cs3Cl profiles are typically identified in
  this manner (though this value rises to \texttildelow3\,--\,5\% in the
  tropics on a typical day), and cloud artifacts may be especially
  prevalent at 147\,hPa.

\item[Quality:] \qualityrule{1.3}

  This threshold for \texttt{Quality} (see
  section~\ref{sec:StatusQuality}) typically excludes less than 1\% of
  CH\cs3Cl profiles on a daily basis; note that it potentially discards
  some ``good'' data points while not necessarily identifying all
  ``bad'' ones.

\item[Convergence:] \convergencerule{1.05}

  On a typical day this threshold for \texttt{Convergence} (see
  Section~\ref{sec:StatusQuality}) discards very few (0.3\% or less) of
  the CH\cs3Cl profiles, most (but not all) of which are also filtered
  out by the other quality control measures.
\end{description}

\subsection*{Artifacts}

\begin{itemize}
\item To be determined.
\end{itemize}

\subsection*{Desired improvements for future data version(s)}

\begin{itemize}
\item To be determined.
\end{itemize}

\begin{table*}[bp]
\caption{Summary of Aura MLS \thisv CH\cs3Cl Characteristics}\vspace{0.5ex}
\label{tab:CH3ClSummary}
\begin{minipage}{\linewidth}
\renewcommand{\thefootnote}{\thempfootnote}
\centering\begin{tabular}{ccccl}
\toprule
\parbox[m]{18mm}{\bfseries\centering{Pressure\\ / hPa}} &
\parbox[m]{24mm}{\bfseries\centering{Resolution\\ V $\times$ H%
\footnote{Vertical and Horizontal resolution in along-track direction.} \\ / km}} &
\parbox[m]{18mm}{\bfseries\centering{Precision%
    \footnote{Precision on individual profiles.} \\/ pptv}} &
\parbox[m]{20mm}{\bfseries\centering{Accuracy\\ / pptv}} &
\parbox[m]{40mm}{\bfseries\centering{Known Artifacts\\ or Other Comments}} \\
%
\midrule
%
3.2\,--\,0.001 & ---   & ---    & ---     & for scientific
  use\\
%
15\,--\,4.6  & 7\,--\,10 $\times$ 550\,--\,850 & $\pm$100 & TBD & \\
%
100\,--\,22  & 4\,--\,6 $\times$ 450\,--\,500 & $\pm$100 & TBD &  \\
%
147      & 4.5 $\times$ 600 & $\pm$100 & TBD & \\
%
1000\,--\,215  & ---    & ---     & ---     & retrieved\\
%
\bottomrule
\end{tabular}
\end{minipage}
\end{table*}

%%% Local Variables: 
%%% mode: latex
%%% TeX-master: "v4-x-report"
%%% End: 

